% !PRINTED 551
\weberpage{551}
Stellen wir, wie in \S.\,163, in der Reihe der Wurzeln
$\alpha, \alpha_1, \ldots \alpha_{m-1}$ jede rational durch die vorangehende dar:
\[
\alpha_1 = \Theta(\alpha), \quad \alpha_2 = \Theta(\alpha_1), \ldots \alpha_{m-1} = \Theta(\alpha_{m-2}), \quad \alpha = \Theta(\alpha_{m-1}),
\]
worin, da hier $\Omega$ reell vorausgesetzt ist, $\Theta(x)$ eine reelle rationale Function von $x$ bedeutet, so folgt zunächst, dass, wenn
eine der Wurzeln reell ist, auch alle übrigen reell sein müssen.
Dies findet immer bei ungeradem $m$ statt, da eine reelle Gleichung ungeraden Grades immer wenigstens eine reelle Wurzel
haben muss.

Bei geradem $m$ können auch imaginäre Wurzeln vorhanden
sein, und wenn eine Wurzel imaginär ist, so müssen es alle sein,
da, wenn eine reelle Wurzel vorkommt, alle anderen auch reell
sind. Bezeichnen wir mit $\Theta^{\nu}(x)$ die $\nu$ malige Wiederholung der
Function $\Theta$, so ist
% [sic] authority prints alpha^k; the later parallel occurrence alpha_k confirms k should be a subscript.
\[
\alpha_{\nu+k} = \Theta^{\nu}(\alpha^k).
\]
Ist also $\alpha = \alpha_0$ mit $\alpha_k$ conjugirt imaginär, so sind auch für
jedes $\nu$ die Functionen $\Theta^{\nu}(\alpha_0)$ und $\Theta^{\nu}(\alpha_k)$, d. h. $\alpha_{\nu}$ und $\alpha_{\nu+k}$
conjugirt imaginär.

Daraus folgt, dass $2k = m$ sein muss, und wir schliessen,
dass im Falle imaginärer Wurzeln
\[
\alpha_k \quad \text{und} \quad \alpha_{k + \frac{m}{2}}
\]
für jeden Index $k$ ein Paar conjugirt imaginärer Wurzeln bilden.

Die cubischen Gleichungen werden durch Adjunction der
Quadratwurzel aus der Discriminante cyklische Gleichungen.
Wenn die Discriminante positiv ist, so sind die Wurzeln in
Uebereinstimmung mit diesem Satze reell.

\begin{center}
\S.\,166.\\
\gesperrt{Theilung des Winkels.}
\end{center}

Zu den cyklischen Gleichungen gehören auch die Gleichungen,
von denen die Theilung eines Winkels in $m$ gleiche Theile abhängt, auf die wir die allgemeinen cyklischen Gleichungen in
einem reellen Körper zurückgeführt haben.

Die Aufgabe kann so formulirt werden:

\gesperrt{Wenn $\cos m\varphi$ und $\sin m\varphi$ gegeben sind, so sollen
daraus $\cos\varphi$ und $\sin\varphi$ gefunden werden.}
